\section{Extending the infinite-padding deficit LC to the entire tail}
This is a further written unbounded refinement of the registered
infinite-padding theorem. It does NOT handle the finite-padding deficit
LC target or solve general C/A2/original UI. Its registration is separate.

Use the notation $Q_u,q$ of
\texttt{binary\_CI\_infinite\_padding\_entrance.tex}.

\begin{lemma}[The final binomial-division kernel row]
For $u\ge3$, $s=u-1$, and EVERY $0\le k\le s$,
\[
 Q_sQ_k-Q_{s+1}Q_{k-1}
       \ge q^s(Q_{k+1}+2qQ_k+q^2Q_{k-1}).             \tag{1}
\]
\end{lemma}
\begin{proof}
First let $u\ge6$. The earlier ULC ratio bound gives
\[
 Q_sQ_k-Q_{s+1}Q_{k-1}\ge\frac12 Q_sQ_k,
\]
since its exact factor is $(u-k)(u+1)/(u(u-k+1))\ge1/2$;
for $k=0$ the factor is one directly.
Also $Q_s=(1-(1-q)^u)/q\ge1+q^2+q^4=4-4q\ge3/2$,
using $q<5/8$, which follows from $q^2+q=1$.
The prior bound $T_k/Q_k\le u+q$ holds up to $k=u-2$.
For $k=u-1$, the ULC upper bound $r_u\le(u-q)/u^2\le1$
and the earlier lower ratio bound $r_{u-1}\ge1/(u-1)$ give
\[
 T_{u-1}/Q_{u-1}\le1+2q+q^2(u-1)\le u+q.
\]
Finally $g_u=q^{u-1}(u+q)$ decreases for $u\ge6$, since
$g_{u+1}/g_u\le(2/3)(7/6)=7/9$. At six,
$g_6\le(2/3)(2/5)^2(20/3)=32/45<3/4$.
Thus the right side of (1) is at most $(3/4)Q_k$, paid by
$(1/2)Q_sQ_k$.
More precisely the displayed estimates give the uniform strict margin
$(7/180)Q_k$ for $u\ge6$.

For the twelve remaining cases, the exact margins in basis $(1,\phi)$
are as follows, in order $k=0,\ldots,u-1$:
\[
\begin{array}{c|l}
u&\text{margins}\\\hline
3&(5,-3),(26,-15),(27,-16)\\
4&(28,-17),(81,-48),(103,-61),(55,-33)\\
5&(38,-23),(264,-160),(574,-349),(588,-359),(277,-170).
\end{array}
\]
Using $\phi<13/8$, their minimum lower bounds are respectively
$1/8,3/8,5/8$, all positive.
\end{proof}

\begin{corollary}[The last robust coefficient centre]
If $L$ is nonnegative LC, $L_0=1$, and $\deg L\le u-1$, then
for $B=LQ_u$ and $s=u-1$,
\[
 B_s^2-B_{s-1}B_{s+1}\ge q^s(B_{s+1}+2qB_s+q^2B_{s-1}).
                                                               \tag{2}
\]
\end{corollary}
\begin{proof}
Retain the first Cauchy--Binet row exactly as in the registered
proof and apply (1) at all $k=s-i$. The formerly missing term is
$L_{s+1}=L_u=0$ by the degree bound. No bound on $L_1$ is needed
at this LAST centre alone.
\end{proof}

\begin{theorem}[Entire infinite-padding geometric deficit LC]
For every CI $L=\prod[t]$ with $t\ge3$, degree $h$, and $m\ge h$,
put $C=L\sum_{k=0}^m(z+z^2)^k$ and $b_n=\phi C_{n-1}-C_n$.
Then $b_n$ is positive throughout $m+1\le n\le\deg C+1$,
zero after that support, and LC at EVERY centre $n\ge m+2$.
Consequently truncating it below $m+1$ produces a globally LC,
nonnegative finite polynomial.
\end{theorem}
\begin{proof}
Positivity follows from the registered $C$ tail LC and its strict
sub-$\phi$ ratio at $m+1$.
First consider the binary polynomial alone. Its deficit is positive
from $m+1$ through $2m+1$, and its LC at $m+2,\ldots,2m$
is exactly the registered robust-budget argument with numerator
$L=1$. The endpoint $2m+1$ is automatic, having next value zero.
Thus its truncation below $m+1$ is globally LC.
Convolve this truncation with $L$. Every triple of the actual full
deficit agrees with it from centre $m+h+2$ onward, so ordinary LC
convolution proves that far tail.

In the remaining band $m+2\le n\le m+h+1$, set $s=n-m-1$.
For $1\le s\le h$ use the registered robust budget at $s\le u-2$,
where $u=m+1$. The ONLY extra case is $m=h$, $s=h=u-1$;
it is paid by (2). In each case the exact geometric-deficit minor
is the robust $B$ minor plus the bounded alternating spectral term,
whose magnitude is at most $q^sT_s$. This proves LC on the entire
remaining band. No finite-padding sequence is substituted here.
\end{proof}

\paragraph{Formal and arithmetic scope.}
\texttt{InfinitePaddingFullDeficit.lean} checks five new scoped
standard-axiom statements: the golden-root rational bound, the final
ULC ratio factor, the ENTIRE unbounded geometric-power estimate,
the quantitative large-row payment, and all twelve small exact
margin inequalities. The quotient-polynomial and Cauchy--Binet
identifications, convolution assembly and support argument are the
written proof, not an end-to-end actual-polynomial Lean theorem.
\texttt{audit\_infinite\_padding\_full\_deficit\_LC.py} independently
checks 7257 final-row budgets, 1200 generic LC last-centre budgets
without an $L_1$ bound, and 359943 entire-tail deficit minors on
2500 actual CI contexts. These finite checks corroborate, rather
than replace, the unbounded argument. Finite short padding, general
C, A2 and the original unamplified conjecture remain OPEN.
